\section{Theory}

\subsection{A population rate curve supplies the nonlinearity}\label{sec:nmm}
A neural mass model describes a population through averaged variables
\cite{ruffini2025rosetta,jansen1995}. In the Jansen--Rit (JR) family, a linear synaptic
filter converts presynaptic firing rate $r$ into postsynaptic potential $u$,
\[
\ddot u+2a\dot u+a^2u=Aa\,r,
\]
with inverse synaptic time $a$ (s$^{-1}$) and PSP gain $A$ (mV); inhibitory synapses use
$b$ and $B$ in the same respective units. Summed postsynaptic potentials give the mean input $v$, which
a static transfer curve converts to firing rate \cite{wilsoncowan1972},
\begin{equation}
\Sig(v)=\frac{2e_0}{1+e^{\,\rho(v_0-v)}} .
\label{eq:sigmoid}
\end{equation}
Here $2e_0$ is the maximum rate, $v_0$ the inflection voltage, and $\rho$ the inverse
voltage scale (s$^{-1}$, mV, and mV$^{-1}$, respectively). Primes denote
derivatives with respect to voltage. Nonzero curvature makes this curve a mixer, while its slope sets
the local gain of the recurrent circuit.

The JR column couples pyramidal, excitatory-interneuron, and inhibitory-interneuron
populations, with field input to the pyramidal sigmoid and local field potential
(LFP) proxy $v=y_1-y_2$, where $y_1$ and $y_2$ are the excitatory and inhibitory
postsynaptic potentials onto the pyramidal population. Its external input rate is
$p$ (s$^{-1}$), and $C$ is the dimensionless connectivity scale.
Figure~\ref{fig:models} locates the field input, nonlinear detector, and
recurrent amplifier in the JR and QIF circuits.
\begin{figure}[t!]
\centering
\includegraphics[width=\textwidth]{fig_models_ti.pdf}
\caption{\textbf{Access, detection, and amplification in the two circuits.}
\emph{(a)} JR: access, coupling $\lambda_0$ times filter $H(\omega)$, converts the field to polarization
$s$, and sigmoid curvature mixes the tones.
\emph{(b)} NMM2 PING: current $I_E=\lambda_I E$ (field $E$, coupling $\lambda_I$) enters the excitatory membrane;
$v_E^2$, the squared mean voltage, mixes its excursions. The inset shows the QIF neurons whose
infinite-population limit gives this exact mean field.
Low-pass symbols mark second-order synapses ($h_e,h_i$ in a; AMPA and GABA
with times $\tau_a,\tau_g$ in b). They suppress carriers more strongly than
the generated beat and transmit the steady shift.
Dashed frames enclose recurrent amplifiers, with small-signal gain
proportional to inverse modal damping $1/\gamma$ near the relevant Hopf.
Blue: excitation; red: inhibition; green: field input.}
\label{fig:models}
\end{figure}


\subsection{From applied field to a squared-envelope beat component}\label{sec:coupling}
We neglect propagation and induction in the extracellular field $E$ (V/m),
but retain finite membrane charging. The effective polarization $s$ (mV)
enters the population sigmoid after cellular access,
\begin{equation}
r=\Sig(v+s),\qquad
E(t)=E_1\cos\omega_1 t+E_2\cos\omega_2 t .
\label{eq:coupling}
\end{equation}
Soma, dendrites, AIS, nodes, and terminals contribute parallel inputs with
different couplings and response times. For nearby carriers centered at
$\omega_c=2\pi f_c$,
\begin{equation}
s\simeq\sum_k \boldsymbol\lambda^k_f(\omega_c)\!\cdot\!\mathbf E,\qquad
\boldsymbol\lambda^k_f(\omega)=\boldsymbol\lambda^k_0\,\lvert H_k(\omega)\rvert,
\label{eq:doors}
\end{equation}
where $\boldsymbol\lambda^k_0$ is compartment $k$'s static polarization vector
(a length) and $H_k$ its access filter \cite{ruffini2025rosetta,galangadea2023}.
This amplitude-gain sum assumes aligned compartment phases; otherwise the
complex transfers must be summed. For one dominant entry point,
$\lambda_f=\lambda_0|H|$ and $\varepsilon_j=\lambda_f(\omega_j)E_j$.
Holding these post-access polarizations fixed isolates detection and gain;
at fixed field they inherit $H$'s frequency dependence.

Fast axonal access motivates $\lambda_f\approx\lambda_0$. At the AIS,
$\Sig(v+s)$ represents a spike-initiation threshold shifted from $v_0$ to
$v_0-s$, requiring no backward transmission to the soma. The conversion from
AIS polarization to $s$ and its response time require calibration. A static
detector additionally requires relaxation faster than the carrier.

Expanding about the field-free voltage $v^*$ (mV)
\cite{ruffini2025ham} gives
$\Sig(v^*+s)\approx\Sig(v^*)+\Sig'(v^*)s+\half\Sig''(v^*)s^2+\dots$.
Quadratic mixing leads when $|s|\ll1/\rho\approx1.8$~mV and
$\Sig''(v^*)\ne0$; larger excursions require the full curve.
The linear term preserves the carriers, while
$2\cos\omega_1t\cos\omega_2t=\cos\Omega t+\cos(\omega_1+\omega_2)t$
generates their difference. Averaging over the carrier removes the fast terms
before applying network gain; synaptic filtering suppresses those terms
downstream of detection. Since $\langle\cos^2\omega_jt\rangle_c=\half$,
\begin{equation}
\big\langle\Sig(v^*+s)\big\rangle_c
=\Sig(v^*)+\tfrac14\Sig''(v^*)\,(\varepsilon_1^2+\varepsilon_2^2)
+\tfrac12\Sig''(v^*)\,\varepsilon_1\varepsilon_2\cos\Omega t+O(\varepsilon^4),
\label{eq:taylor}
\end{equation}
giving a steady source from the summed squares and a signed beat-source amplitude (s$^{-1}$ )
\begin{equation}
Y_\Omega\simeq\half\,\Sig''(v^*)\,\varepsilon_1\varepsilon_2 .
\label{eq:demod}
\end{equation}
Here $\varepsilon$ denotes a common small-amplitude scale and $\simeq$
denotes leading order in polarization throughout.

For a dynamic cellular response with synaptic input held fixed, define
the second-order coefficient $\chi^{(2)}$ through the complex source amplitude,
\begin{equation}
Y_\Omega\simeq\chi^{(2)}(\omega_2,-\omega_1)\,
\varepsilon_2\varepsilon_1,\qquad
r_\Omega(t)=\operatorname{Re}\!\left[Y_\Omega e^{i\Omega t}\right].
\label{eq:chi2}
\end{equation}
The generated rate $r_\Omega$ precedes recurrence, so $\chi^{(2)}$ has units
s$^{-1}$mV$^{-2}$. With this cosine-amplitude convention, its memoryless
limit is $\Sig''/2$. Cellular history, state, and response times can make
it complex and frequency-dependent; symmetry or cancellation can remove
second-order mixing, leaving higher orders.

The static detector predicts quadratic scaling, a vanishing source at the
inflection, and carrier independence at fixed polarization. At fixed field,
access attenuation raises the field required for a given source, consistent
with carrier-dependent thresholds \cite{opancar2025,peripheral2023}.

\subsection{Envelope oscillations and spatial weighting}\label{sec:envelope}
Envelope following and quadratic detection assign different spatial doses
to the same field. At position $\mathbf x$, let $\varepsilon_j(\mathbf x)$ be
the signed polarization of tone $j$ along a fixed direction. The nonnegative
envelope of the two-tone polarization is
\[
\mathcal A(\mathbf x,t)=\sqrt{\varepsilon_1(\mathbf x)^2+\varepsilon_2(\mathbf x)^2
+2\varepsilon_1(\mathbf x)\varepsilon_2(\mathbf x)\cos\Omega t}.
\]
At each location, low-pass filtering (LPF) removes the carrier terms and gives
\begin{equation}
\mathrm{LPF}[s^2]=\tfrac12\mathcal A^2
=\frac{\varepsilon_1^2+\varepsilon_2^2}{2}+\varepsilon_1\varepsilon_2\cos\Omega t,
\label{eq:sqenv}
\end{equation}
with spatial arguments suppressed. The squared envelope is sinusoidal at
$\Omega$; the envelope itself is not. For matched tones it is the rectified
cosine $\mathcal A(t)=2\varepsilon|\cos(\Omega t/2)|$.
The conventional TI map measures its excursion,
$\mathcal M[y]\equiv\max_t y-\min_t y$ \cite{grossman2017}.

Table~\ref{tab:transduction-maps} separates these geometric quantities from
detector gains; overbars denote beat-period means. An ideal envelope follower
outputs $c_1\mathcal A$ for gain $c_1$, whereas a smooth transfer $g$ with
nonzero curvature gives $g''(v^*)\mathcal A^2/4$ at weak input ($g=\Sig$ in JR).
Thus the nonnegative mean-square map $\mathcal I$ specifies neither excitation
nor suppression, and an envelope-excursion map alone specifies no mean output.
Envelope following is itself a response assumption.

\begin{table}[ht]
\centering\small
\renewcommand{\arraystretch}{1.4}
\begin{tabular}{@{}>{\raggedright\arraybackslash}p{0.22\textwidth}
 >{\raggedright\arraybackslash}p{0.34\textwidth}
 >{\raggedright\arraybackslash}p{0.36\textwidth}@{}}
\toprule
 & Envelope following & Quadratic detection \\
\midrule
Slow signal & $\mathcal A(\mathbf x,t)$ & $\tfrac12\mathcal A(\mathbf x,t)^2$ \\
Mean before gain & $\overline{\mathcal A}(\mathbf x)$
 & $\mathcal I(\mathbf x)=\dfrac{\varepsilon_1^2+\varepsilon_2^2}{2}$ \\
Modulation excursion & $2\min(|\varepsilon_1|,|\varepsilon_2|)$
 & $2|\varepsilon_1\varepsilon_2|$ \\
Units before detector gain & mV & mV$^2$ \\
Signed steady source (output units) & $c_1\overline{\mathcal A}$
 & $Y_{\rm DC}=\tfrac12g''(v^*)\mathcal I$ \\
\bottomrule
\end{tabular}
\caption{Envelope and quadratic maps at the same position $\mathbf x$.
The first four rows give geometric quantities before detector gain; the last
gives the steady source after gain, before recurrent feedback. Its sign depends
on the detector. Averaging the slow signal over a beat period gives its mean.
Excursion means maximum minus minimum, not Fourier amplitude.}
\label{tab:transduction-maps}
\end{table}

The excursion maps in Table~\ref{tab:transduction-maps} are related pointwise by
\begin{equation}
\mathcal{M}\big[\tfrac12\mathcal A^2\big]=\max(|\varepsilon_1|,|\varepsilon_2|)\;\mathcal{M}[\mathcal A].
\label{eq:reweight}
\end{equation}
The dominant polarization reweights the envelope map. Both maps can have a
central maximum, although the quadratic example in Fig.~\ref{fig:spatial_maps}a
is broader. A larger superficial dominant carrier degrades deep contrast:
for $(\varepsilon_1,\varepsilon_2)=(10,1)$ superficially and $(0.5,0.5)$
at depth, the superficial/deep excursion ratios are $2$ for the envelope
and $40$ for the quadratic map. A sufficiently weak second carrier can
instead suppress the superficial product. Actual contrast requires the
vector fields.

With fixed electrode pairs, frequencies, neural state, linear access, and
projection directions, write
$\varepsilon_j(\mathbf x)=I_j^{\rm pair}\mathcal P_j(\mathbf x)$,
where $I_j^{\rm pair}>0$ is peak pair current (mA) and $\mathcal P_j$ is
signed polarization per current (mV/mA). Then
\[
\mathcal M[\tfrac12\mathcal A^2]
=2I_1^{\rm pair}I_2^{\rm pair}|\mathcal P_1\mathcal P_2|.
\]
Changing the ratio therefore rescales the beat map without moving its maxima
or changing normalized maps and target/reference ratios
\cite{cassara2025part2}. Fixed detector coefficients and linear propagation
preserve the common current factor, even with dynamic mixing and heterogeneous
gains. Steering requires a changed carrier field shape, for example through
current redistribution among electrodes at one frequency.

The steady source provides a way for current-ratio changes to reshape the
beat response: it weights $\mathcal P_1^2$ and $\mathcal P_2^2$ by the squared
currents and can thereby shift the operating points. Such state changes fall
outside the fixed-state argument above. Threshold effects also depart from
the weak quadratic regime; the $2\min$ steering law requires envelope following.

The product law permits an interior beat maximum. Its spatial weighting also
depends on local detector and network factors (Eq.~\eqref{eq:receiver}), so
resonance can compensate for an unfavorable product only where the source
is nonzero.

The scalar maps use a common fixed projection,
$\varepsilon_j=\boldsymbol\lambda\cdot\mathbf E_j$; the beat changes phase
when either projection changes sign. The current-product scaling also holds
when a detector combines input from several compartments. Let $a_j^k$ be the
complex polarization of carrier $j$ in compartment $k$. A weak dynamic
detector has a bilinear beat source
$Y_\Omega=\sum_{kl}B_{kl}(\omega_2,-\omega_1)a_2^k\overline{a_1^l}$,
where $B_{kl}$ are fixed response coefficients and the overbar means complex
conjugation. A single scalar input is a rank-one specialization. Linear
access gives $a_j^k=I_j^{\rm pair}b_j^k$, with $b_j^k$ the polarization per
unit current, so the common factor $I_1^{\rm pair}I_2^{\rm pair}$ and its
fixed spatial pattern survive without a universal projection vector.

\subsection{Detector laws and harmonic signatures}\label{sec:transduction}
For a narrowband input $s(t)=\mathcal A(t)\cos\omega_c t$, a memoryless
nonlinearity $g$ has low-frequency output given by its carrier-cycle average,
\[
Y_{\rm LF}(\mathcal A)=(2\pi)^{-1}\!\int_0^{2\pi}\!g(v^*+\mathcal A\cos\theta)\,\mathrm{d}\theta,
\]
where $\theta$ is carrier phase. This transduction contains
only even powers,
\begin{equation}
Y_{\rm LF}(\mathcal A)=g(v^*)+\tfrac14 g''(v^*)\mathcal A^2+\tfrac1{64}g^{(4)}(v^*)\mathcal A^4+\cdots
\label{eq:cycavg}
\end{equation}
Thus smooth detection with nonzero curvature is quadratic at weak input.
A threshold detector instead gives
$\pi^{-1}\arccos(\delta/\mathcal A)$ for envelope amplitude
$\mathcal A>\delta$ and zero below, where $\delta>0$ is its offset.
Other rectifier laws are compared in the supplement (\S\ref{supp:harmonics}).

At matched tones, envelope following is linear in amplitude with
second/fundamental harmonic ratio $1/5$. Quadratic detection gives only
the steady term and fundamental; its quartic correction gives
$Y_{2\Omega}/Y_\Omega\sim(c_4/c_2)\varepsilon^2$, with
$c_2=g''/4$ and $c_4=g^{(4)}/64$. A threshold detector is silent for
$2\varepsilon\le\delta$ and has an amplitude-dependent ratio above onset.
The full harmonic series and mixed-detector inversion are supplementary
(\S\ref{supp:harmonics}).
\ifjne\else
\subsection{Harmonic ladder and mixed detectors}\label{supp:harmonics}
Full- and half-wave rectifiers yield $2\mathcal A/\pi$ and
$\mathcal A/\pi$, respectively; a diode-and-capacitor detector follows the
peak. Channel rectification, threshold crossing, and release events are
possible biological implementations.

Envelope following and square-law detection are distinguishable at a single site.
For equal polarization amplitudes, $\varepsilon_1=\varepsilon_2=\varepsilon>0$, $\mathcal A(t)=2\varepsilon|\cos(\Omega t/2)|$.
The Fourier series
\[
|\cos\theta|=\frac{2}{\pi}+\frac{4}{\pi}\sum_{n\ge1}\frac{(-1)^{n+1}}{4n^2-1}\cos(2n\theta)
\]
gives nonnegative harmonic amplitudes $A_{n\Omega}$ with $A_{2\Omega}/A_\Omega=1/5$, $A_{3\Omega}/A_\Omega=3/35$ and
$A_{4\Omega}/A_\Omega=1/21$. The square-law output $\half \mathcal A^2=\varepsilon^2(1+\cos\Omega
t)$ contains only DC and $\Omega$. For the mixed family $T(\mathcal A)=c_1\mathcal A+c_2\mathcal A^2$, with
$c_1\ne0$ and a nonzero fundamental, the
pre-network ratio is
\begin{equation}
\mathcal H_{\rm det}\equiv\left|\frac{Y_{2\Omega}}{Y_\Omega}\right|
=\frac{1/5}{\left|1+\tfrac{3\pi}{4}\,c_2\varepsilon/c_1\right|},
\label{eq:rdet}
\end{equation}
so measuring $\mathcal H_{\rm det}$ constrains $c_2\varepsilon/c_1$ within this family.\footnote{The inverse has two branches unless the sign of the fundamental
relative to the linear term is known; for $c_2/c_1\ge0$ it is unique.}

\fi

The network modifies these ratios through
$\mathcal H_{\rm meas}=\mathcal H_{\rm det}|G(2\Omega)/G(\Omega)|$.
Near resonance it can suppress an envelope follower's second harmonic,
making it appear quadratic. Amplitude sweeps and gain correction are
therefore needed; a threshold detector can match one harmonic ratio, so that
ratio alone cannot identify the detector.

\subsection{The steady source has a state-dependent sign}\label{sec:dc}
Detection converts nonnegative mean-square polarization into a signed steady
source; recurrence converts that source into a mean response. ``DC-like''
describes this source, not the applied field, which has no DC component.
At matched polarizations $\varepsilon>0$,
the envelope and quadratic means are $4\varepsilon/\pi$ and
$\varepsilon^2$ before detector gains (supplement, \S\ref{supp:envelope-proof}). Either ideal detector gives a mean
output with one carrier, but no beat; adaptation can alter the mean response.

The static JR source is
\[
Y_{\rm DC}(\mathbf x,\mathbf n)
=\frac{\Sig''(v^*_{\mathbf x})}{2}\,\mathcal I(\mathbf x,\mathbf n),
\qquad
\mathcal I(\mathbf x,\mathbf n)
=\frac{\varepsilon_1(\mathbf x,\mathbf n)^2+
\varepsilon_2(\mathbf x,\mathbf n)^2}{2}\ge0,
\]
where $v^*_{\mathbf x}$ is the local baseline and $\mathbf n$ the common
unit cell orientation. Convex and concave rate curves give positive and
negative sources, respectively; the quadratic source vanishes at the inflection
and in the saturated tails.
Reversing either carrier or $\mathbf n$ leaves the squares unchanged;
other orientations change magnitudes but not the sign at fixed curvature.
Figure~\ref{fig:spatial_maps}c shows both signs on one positive normalization.
At the pinned JR state, $v^*=8.30$~mV, the source lowers the local rate
because $\Sig''(v^*)<0$.
This curvature rule applies to static detection; dynamic rectification and
recurrence can change the sign of the mean response.
The same operating-point dependence enters the multitone sigmoid expansion for
endogenous infraslow ($<0.1$~Hz) mixing in coupled Jansen--Rit models \cite{Just2026mixing}.

\label{sec:opp}
The Population Excitation--Inhibition Index (PEIX) \cite{ruffini2025cfc}
measures signed departure from maximal sigmoid slope, positive above the
inflection and negative below; its definition and curvature identities are
in Appendix~\ref{app:peix}. Changes in $v^*$ through neuromodulation,
attention, arousal, psychedelic exposure, or dysfunction of PV interneurons
in Alzheimer's disease can scale or reverse the static JR source.
Dynamic detectors have their own state-dependent coefficients and susceptibilities, without
a universal PEIX sign rule; state sweeps isolate detection only if recurrent
gain is also measured.

Carrier averaging retains the beat cross term; a beat-period average removes
it, leaving intensities without cross-frequency cancellation. Where
$Y_{\rm DC}\ne0$,
$|Y_\Omega/Y_{\rm DC}|=2|\varepsilon_1\varepsilon_2|/
(\varepsilon_1^2+\varepsilon_2^2)\le1$, with equality at matched magnitudes.
A common access factor cancels. In the superficial/deep example of \S\ref{sec:envelope}, $\mathcal I$ is
$50.5$ superficially and $0.25$ at depth, a $202$-fold contrast;
the corresponding beat/DC ratios are $0.20$ and $1$.

In a homogeneous, isotropic, source-free conductor, the potentials $\phi_j$,
field components $\mathbf E_j=-\nabla\phi_j$, and polarizations
$\varepsilon_j=\boldsymbol\lambda\cdot\mathbf E_j$ for spatially fixed
coupling $\boldsymbol\lambda$ are harmonic. Consequently
$\nabla^2\mathcal I=\sum_j|\nabla\varepsilon_j|^2\ge0$:
the intensity is subharmonic, with no strict interior maximum. Its maximum
lies on the region's boundary, including conductivity interfaces, regardless
of montage. With spatially constant nonzero detection, the same holds for
$|Y_{\rm DC}|$; negative curvature makes the boundary source most negative.
Spatially varying access, orientation, or detector coefficients require a
weighted map, and the restriction does not extend to the physiological mean
response or an indefinite multicompartment DC form.

The envelope mean is also subharmonic: at each beat phase it is a norm of
harmonic components, and phase averaging preserves that property (supplement,
\S\ref{supp:envelope-proof}).
\ifjne\else
\subsection{Envelope-mean subharmonicity}\label{supp:envelope-proof}
The envelope mean is
\[
\overline{\mathcal A}(\mathbf x)=\frac{1}{2\pi}\int_0^{2\pi}
\sqrt{\varepsilon_1(\mathbf x)^2+\varepsilon_2(\mathbf x)^2
+2\varepsilon_1(\mathbf x)\varepsilon_2(\mathbf x)\cos\vartheta}\,\mathrm d\vartheta,
\]
where $\vartheta$ is beat phase.

For fixed scalar coupling in homogeneous, isotropic, source-free tissue,
the envelope mean is subharmonic. At each beat phase $\vartheta$,
$\mathcal A$ is the Euclidean norm of
$(\varepsilon_1+\varepsilon_2\cos\vartheta,\,\varepsilon_2\sin\vartheta)$,
a vector of harmonic functions. A norm is convex, so this envelope is subharmonic;
averaging over $\vartheta$ preserves that property. Thus both idealized mean maps
lack a strict interior maximum under these conditions (Fig.~\ref{fig:spatial_maps}b).

\paragraph{Interior maxima of the excursion and the beat amplitude.}\label{supp:interior-maxima}
The product of the polarizations obeys
$\nabla^2(\varepsilon_1\varepsilon_2)=2\nabla\varepsilon_1\cdot\nabla\varepsilon_2$,
whose sign is unrestricted, and an excursion, a temporal maximum minus a temporal
minimum, need not be subharmonic. A local example shows that both can have a strict
three-dimensional interior maximum under the conditions above. In dimensionless
coordinates $\mathbf x=(x,y,z)$, take $0<c<1/2$, $q=x^2-\tfrac12(y^2+z^2)$, and the harmonic potentials
\[
\phi_{1,2}=-x\mp\tfrac12(x^2-y^2)-c\left[\tfrac13x^3-\tfrac12x(y^2+z^2)\right].
\]
With unit coupling along the fixed $x$ direction,
$\varepsilon_{1,2}=-\partial_x\phi_{1,2}=1\pm x+cq$, positive near the origin. Then
\[
\varepsilon_1\varepsilon_2=1+(2c-1)x^2-c(y^2+z^2)+O(|\mathbf x|^4),
\]
with a negative-definite Hessian at the origin, and the envelope excursion
$2\min(\varepsilon_1,\varepsilon_2)=2(1+cq-|x|)$ is also strictly maximal there.
This is an existence example, not a prediction for a montage. The subharmonic maps
above can still have interior minima; the restriction concerns maxima only.

\fi
By contrast,
$\nabla^2(\varepsilon_1\varepsilon_2)=2\nabla\varepsilon_1\cdot\nabla\varepsilon_2$
has unrestricted sign. Within this homogeneous scalar model, interior
selectivity must come from the oscillatory response, not either mean map.

With isotropic orientations and orientation-independent access and curvature,
orthogonal fields cancel the averaged beat while retaining steady intensity
(supplement, \S\ref{supp:orientation}).
\ifjne\else
\subsection{Isotropic orientation averages}\label{supp:orientation}
For isotropically distributed unit orientations
$\mathbf n$ with orientation-independent access and curvature,
$\langle(\mathbf n\cdot\mathbf E_j)^2\rangle=|\mathbf E_j|^2/3$
and $\langle(\mathbf n\cdot\mathbf E_1)(\mathbf n\cdot\mathbf E_2)\rangle
=\mathbf E_1\cdot\mathbf E_2/3$: orthogonal fields cancel the averaged beat,
but retain the DC intensity.

\fi

An unmodulated carrier at the same local RMS polarization gives the same
second-order mean shift if its DC mixing coefficient equals those of both
TI carriers. This motivates the equal-RMS carrier control in
\S\ref{sec:qif-methods}. For a dynamic detector, the coefficient must vary
negligibly over the carrier interval.

\subsection{Comparison with direct tDCS and tACS}\label{sec:direct-comparison}
We compare TI's steady and beat sources with transcranial direct-current
stimulation (tDCS) and tACS at $\Omega$, respectively, using polarizations
at the same sigmoid entry point. For equal TI tones
$\varepsilon_1=\varepsilon_2=\varepsilon>0$,
\[
Y_{\rm DC}^{\rm TI}\simeq Y_\Omega^{\rm TI}
\simeq \frac12\Sig''(v^*)\varepsilon^2.
\]
A constant polarization $\varepsilon_{\rm D}$ or sinusoid
$\varepsilon_{\rm A}\cos\Omega t$ instead gives
\[
Y_{\rm DC}^{\rm tDCS}\simeq\Sig'(v^*)\varepsilon_{\rm D},
\qquad
Y_\Omega^{\rm tACS}\simeq\Sig'(v^*)\varepsilon_{\rm A}.
\]
Here $\varepsilon_{\rm D}$ is signed and negative $\varepsilon_{\rm A}$
denotes phase reversal. The equivalent direct polarization matching either
TI source is
\[
\varepsilon_{\rm eq}
=\frac{\Sig''(v^*)}{2\Sig'(v^*)}\,\varepsilon^2,
\qquad
\frac{\Sig''(v^*)}{\Sig'(v^*)}
=-\rho\tanh\!\left[\frac{\rho(v^*-v_0)}2\right].
\]
Sigmoid steepness and operating voltage determine the ratio; the rate ceiling
cancels, and the quadratic TI sources vanish at the inflection.

For the JR parameters $\rho=0.56$~mV$^{-1}$, $v_0=6$~mV,
$e_0=2.5$~s$^{-1}$, and the pinned operating point $v^*=8.30$~mV,
$\Sig'=0.4745$~s$^{-1}$mV$^{-1}$ and
$\Sig''=-0.1508$~s$^{-1}$mV$^{-2}$. With each direct-input amplitude
equal to the per-carrier TI amplitude, the source-magnitude ratio is
$|\varepsilon_{\rm eq}|/\varepsilon
=0.1589\,\varepsilon/(1~\mathrm{mV})$:
$1.59\%$ at $0.1$~mV and $4.77\%$ at $0.3$~mV, rounded to $2$--$5\%$.
At $0.3$~mV per carrier, the equivalent direct polarization is
$-0.0143$~mV, or $-14.3~\mu$V. The negative sign specifies a suppressive
local source for tDCS and an inverted source phase for tACS.

Matching the total TI peak instead uses direct amplitude $2\varepsilon$
and halves these ratios. RMS matching uses tACS amplitude
$\sqrt2\varepsilon$ or tDCS magnitude $\varepsilon$, because TI's RMS
polarization is $\varepsilon$. The tACS quadratic steady source,
$\Sig''\varepsilon_{\rm A}^2/4$, then equals TI's for this static detector.
Equal electrode currents need not match polarization: access differs among
DC, beat and carrier frequencies, and among entry compartments.

At a common state and entry point, matched steady sources share the
zero-frequency transfer, and matched beat and tACS sources share the transfer
at $\Omega$. These transfers differ, and rectification of the amplified TI
beat can add a mean shift.
Source ratios therefore imply neither universal physiological-effect ratios
nor equal spatial responses.

\subsection{Amplification: the synaptic low-pass and the network resonance}\label{sec:carrier}
Synaptic filtering suppresses the carrier after the nonlinearity has generated
the beat. Each JR excitatory synapse has transfer function
$H_s(\omega)=Aa/(i\omega+a)^2$, a second-order low-pass with characteristic
frequency $a/(2\pi)\approx16$~Hz. It passes a $10$~Hz component with relative
amplitude about $0.7$, while reducing a $100$~Hz component to about $2.5\%$. Provided the carrier does not itself excite a responsive
mode, the loop therefore reads predominantly the generated low-frequency drive.

% fig:carrier moved to Results (Theory is equations-only)

The recurrent circuit selects the recovered beat. Linearizing about the JR
fixed point gives a leading pair $\mu_\pm=-\gamma\pm i\omega_0$, with
modal damping $\gamma$ and frequency $\omega_0$. Below a supercritical
Hopf, $\gamma\to0^+$ and the quality factor
$\mathcal Q=\omega_0/(2\gamma)$ grows. A constant-numerator two-pole
approximation to the complex susceptibility $G$ gives
\begin{equation}
|G(\Omega)|\propto\big[(\omega_0^2+\gamma^2-\Omega^2)^2+(2\gamma\Omega)^2\big]^{-1/2},
\qquad |G(\omega_0)|\propto\frac{1}{2\gamma\omega_0}
\quad(\gamma\ll\omega_0),
\label{eq:gain}
\end{equation}
an approximately Lorentzian power peak near $\omega_0$ with $1/\gamma$
growth at a simple Hopf with nonzero input--output projections.\footnote{
The eigenvalue frequency $\omega_0$ differs from the canonical oscillator's
undamped frequency $\sqrt{\omega_0^2+\gamma^2}$ and peak
$\sqrt{\omega_0^2-\gamma^2}$ ($\omega_0>\gamma$).
Frequency-dependent projections can further shift the peak.}
This weak-drive law requires stability after the rectified mean shift.
Above the Hopf, entrainment instead means phase capture of an autonomous rhythm.

Write $D$ for the detector's beat-source coefficient: for the JR rate source,
$D=\half\Sig''(v^*)$, half the local sigmoid curvature. Combining access,
detection, and amplification gives
\begin{equation}
R_\Omega\simeq G(\Omega;\gamma)\;D\;\lambda_f(\omega_1)E_1\,\lambda_f(\omega_2)E_2 ,
\label{eq:receiver}
\end{equation}
% Retain the former spatial-equation label as an alias of the general receiver.
\makeatletter\begingroup\edef\@currentlabel{\theequation}\label{eq:spatial}\endgroup\makeatother
where $\lambda_f$ maps field to local carrier polarization, $D$ converts the
product of those polarizations into a beat drive, and $G$ maps that drive to the
complex observed response, whose measured amplitude is $A_\Omega=|R_\Omega|$.
For JR, $G$ is the response to a generated rate drive.
A static detector has carrier-independent $D$ at fixed polarization;
a dynamic detector can retain frequency dependence. A carrier sweep and a beat sweep
therefore constrain different stages rather than a single effective gain.

Resonance requires a beat source. The cubic Hopf normal form
responds to an in-band drive but has no second-order difference-frequency source. A large resonant susceptibility alone therefore does not
establish beat generation.
In the calibrated noise-free Stuart--Landau networks tested in \cite{Ruffini2026iso},
first-order rate transmission and readout yielded negligible infraslow BOLD power.

\subsection{Carrier access depends on the effective membrane response}\label{sec:khz}
Carrier independence at fixed polarization does not imply independence at fixed
applied field. As an isopotential access model, consider
$H_m(\omega)=1/(1+i\omega\tau_m)$, with effective time constant $\tau_m$.
This is the response of a point membrane to injected current; a field coupled
through a cable need not have the same transfer.
For nearby carriers, the two polarizations acquire nearly the same attenuation.
The quadratic beat then satisfies
\begin{equation}
\frac{|Y_\Omega(f_c)|}{|Y_\Omega^{\rm direct}|}
\simeq |H_m(\omega_c)|^2
=\frac{1}{1+(\omega_c\tau_m)^2}
\sim\frac{1}{(\omega_c\tau_m)^2}
\quad(\omega_c\tau_m\gg1),
\label{eq:khz}
\end{equation}
where $Y_\Omega^{\rm direct}$ is the beat generated with the same static coupling
and no access attenuation. This filter belongs either in the field-to-polarization
map or in the membrane equation that computes that polarization.

For an illustrative somatic $\tau_m=16$~ms, $|H_m|^2$ is about
$10^{-4}$, $2.5\times10^{-5}$, and $10^{-6}$ at $1$, $2$, and $10$~kHz.
The assumed $0.2$~ms case gives $0.39$ and $0.14$ at $1$ and $2$~kHz
(Fig.~\ref{fig:khz}); this transfer has not been measured under TI at an
AIS or node. Table~\ref{tab:doors} distinguishes measured responses,
assumed access, and downstream filtering. Axonal decay alone specifies
neither field coupling nor a TI cutoff.

A sufficiently long passive cable with a sealed end in a uniform axial field
has quadratic-beat factors near $10^{-2}$ at $1$~kHz and $10^{-3}$ at
$10$~kHz for $\tau_m=16$~ms, exceeding the point-membrane values.

QIF computes access in its membrane equation, requiring no additional filter
on current input. Shorter access improves transmission without determining
mixing or recurrent gain.

\begin{table}[h]
\centering\small
\renewcommand{\arraystretch}{1.15}
\begin{tabular}{@{}>{\raggedright\arraybackslash}p{2.8cm}
>{\raggedright\arraybackslash}p{3.5cm}
>{\raggedright\arraybackslash}p{8.6cm}@{}}
\toprule
Response & Time scale & Evidence and role \\
\midrule
Somatic point-membrane model & $\tau_m=16$~ms; $f_{\rm char}\approx10$~Hz &
Illustrative current-driven filter, informed by somatic measurements near $12$ and $16.5$~ms
\cite{deitcher2017,eyal2016}. \\
Myelinated axon & minimum fitted decay $\approx1$~ms &
Injected-current responses in Cohen et al.'s periaxonal circuit
\cite{cohen2020periaxonal}; evidence for distinct axonal voltage dynamics. \\
Fast axonal access & assumed $\tau_m=0.2$~ms; $f_{\rm char}\approx800$~Hz &
Reduced-model sensitivity case for a fast compartment, such as an AIS or node;
TI coupling and charging time require calibration. \\
Prior TI network model & $\tau_m=1$ and $0.4$~ms &
Short effective membrane times governing kHz sensitivity in a gamma network
\cite{esmaeilpour2021}; model parameters, not measured axonal TI decays. \\
Field-induced charging & $\approx0.7~\mu$s &
Measured with strong microsecond field pulses in isolated neurons
\cite{semenov2025}; a different stimulation regime. \\
Cellular detector & channel- and state-dependent &
Luff et al.'s TTX-sensitive mixing persists near $5$~kHz
\cite{luff2024mixer}; this measures detection, not a passive time constant. \\
JR excitatory synapse & $\tau_s=10$~ms; $f_{\rm char}\approx16$~Hz &
Model filter acting on the generated rate, downstream of detection. \\
\bottomrule
\end{tabular}
\caption{\textbf{Access, detection, and downstream filtering involve different response
times.} Measured decays, pulse-charging times, and assumed model parameters are
identified separately. The characteristic frequency $f_{\rm char}=1/(2\pi\tau)$
is given only for the specified filters; an axonal decay is not automatically a TI
cutoff. AIS, axon initial segment; TTX, tetrodotoxin.}
\label{tab:doors}
\end{table}

Any additional gate filter must represent kinetics absent from the detector
and have a model or measurement behind it. QIF lacks explicit sodium-channel
kinetics; neither it nor a static sigmoid specifies a universal $1$~kHz cutoff.

\subsection{The exact mean field extends the static detector}\label{sec:dynamic}
The exact mean field of all-to-all quadratic integrate-and-fire (QIF) neurons
with Lorentzian excitabilities supplies both a static rate curve and its
dynamic extension. For weak carriers at fixed synaptic input, steady
rectification follows the curve's curvature at low frequency and its slope
at high frequency. The Montbri\'o--Paz\'o--Roxin (MPR) equations evolve physical
rate $r$ and dimensionless mean voltage $v$ \cite{montbrio2015}; dynamic
synapses produce a second-generation neural mass (NMM2)
\cite{ruffini2021nmm2,clusella2023}. The population equations are
\begin{equation}
\tau_m\dot r=\frac{\Delta}{\pi\tau_m}+2rv,\qquad
\tau_m\dot v=v^2+\bar\eta-(\pi\tau_m r)^2+I.
\label{eq:mpr}
\end{equation}
Here $\tau_m$ is membrane time, dots denote physical-time derivatives,
$\bar\eta$ and $\Delta>0$ are the excitability center and half-width, and
$I$ is dimensionless input, including recurrent drive when present.
For constant total drive $Q=\bar\eta+I$, equilibrium gives
$v=-\Delta/(2\pi\tau_m r)$ and
\begin{equation}
\tau_m r=F(Q)=\frac{1}{\pi}\sqrt{\frac{Q+\sqrt{Q^2+\Delta^2}}{2}} .
\label{eq:Fmpr}
\end{equation}
Primes on $F$ denote drive derivatives. Below its inflection at
$Q=\Delta/\sqrt3$, curvature reflects recruitment across threshold;
above it, the concave square-root rate of active cells. This nonlinearity
exists without recurrence and grows without bound, unlike the JR sigmoid.

If relaxation is faster than synaptic variation, NMM2 reduces to
$r\simeq F(Q)/\tau_m$ with dynamically filtered recurrence
\cite{ruffini2021nmm2,clusella2023}. For population relaxation time
$\tau_{\rm rel}$ and synaptic time $\tau_s$ (Eq.~\eqref{eq:slow_synapse},
Appendix~\ref{app:models}), this requires $\tau_{\rm rel}\ll\tau_s$.
Applying the static curve to carriers further requires
$\omega_c\tau_{\rm rel}\ll1$: the same population can track synaptic
input quasistatically while responding dynamically to kHz input.
Recurrent excitatory--inhibitory circuits resonate without intrinsic
cellular resonance \cite{brunelhakim1999,brunelwang2003}. Such synaptic
modes survive the fast-membrane limit, whereas modes requiring finite
membrane dynamics need not.

\paragraph{Field coupling and recurrent gain in the slow limit.}
The field input survives this reduction. For one self-coupled population,
write $I=I_b+\tau_mJr_{\rm syn}+\lambda_I E$, where $I_b$ is the normalized
background input, $r_{\rm syn}$ the filtered rate, and $J$ dimensionless
self-coupling. For a fixed field direction, $\lambda_I=\kappa_E/I_s$:
$\kappa_E$ is equivalent current per field and $I_s$ the current unit.
A cellular voltage-polarization coefficient requires an impedance conversion
before entering this current-driven model. At equilibrium $r_{\rm syn}=r$;
differentiating $\tau_mr=F(\bar\eta+I_b+\tau_mJr+\lambda_IE)$ gives
\begin{equation}
\frac{\partial r_*}{\partial E}
=\frac{\lambda_I F'(Q_*)}{\tau_m[1-JF'(Q_*)]}.
\label{eq:slow-field-gain}
\end{equation}
Positive self-feedback amplifies the steady response on the stable branch
$JF'(Q_*)<1$; negative feedback suppresses it. The bare curve $F$ takes
total drive as its argument. A transfer fitted after an internally recurrent
population settles already contains that feedback and must not be coupled
through the same loop again. A JR-like description retains this gain when
input conversion, local slope, and feedback are matched
\cite{ruffini2021nmm2,clusella2023}; its connections between populations do
not automatically specify the same within-population self-coupling.

Outside the adiabatic limit, the quadratic voltage term mixes the carriers
before the rate responds. For an uncoupled population under a weak current
$I_{\rm ac}\cos\omega t$, write the mean-rate shift as
$I_{\rm ac}^2\kappa_{\rm DC}/\tau_m$. At slow carriers,
$\kappa_{\rm DC}\to F''(Q)/4$; at fast carriers,
$\kappa_{\rm DC}\sim F'(Q)/[2(\omega\tau_m)^2]$.
The slow response follows curvature and can be negative above the inflection.
The fast DC coefficient instead follows the positive slope $F'$, largest at the
inflection, without a PEIX-based sign reversal: voltage excursions of order
$I_{\rm ac}/(\omega\tau_m)$ generate a positive input shift through $v^2$.
A shorter membrane time can therefore change both the magnitude and sign
of the mean response at a given physical carrier frequency.

This fast-carrier attenuation includes membrane access. Normalizing by the
squared voltage excursion leaves a coefficient tending to $F'/2$.
For two nearby current tones of amplitudes $I_j$, voltage squaring gives a
beat source proportional to $I_1I_2/(\omega_1\omega_2\tau_m^2)$.
The dynamic rate equation further attenuates the carriers after mixing;
that carrier attenuation is not applied to the generated beat. The crossover is derived in
Appendix~\ref{app:mpr-coefficients}; fixed-voltage coefficients and carrier
exponents are supplementary (\S\ref{supp:dynamic-details}).

In this cellular realization of quadratic detection, the unit coefficient of $v^2$ in membrane-time units
plays the role of $\Sig''/2$; $G$ includes the voltage and rate dynamics
converting that source into a rate beat. Other smooth cellular mechanisms
share the weak-input product law without sharing its coefficient, phase,
or carrier dependence.

